\documentclass[11pt,letterpaper]{article}
\usepackage[utf8]{inputenc}
\usepackage[T1]{fontenc}
\usepackage{lmodern}
\usepackage{amsmath}
\usepackage[spanish,es-tabla]{babel}
\usepackage[margin=2.5cm]{geometry}
\usepackage{microtype}
\usepackage{booktabs,longtable,array}
\usepackage[skip=0.5em]{parskip}
\usepackage{xurl}
\usepackage[hidelinks]{hyperref}
\usepackage{tikz}
\usetikzlibrary{positioning,arrows.meta}
\usepackage{listings}
\lstset{basicstyle=\ttfamily\small,breaklines=true,columns=fullflexible,
  keepspaces=true,showstringspaces=false,frame=single}
\newcommand{\archivo}[1]{\path{#1}}
\title{SmartCity Analytics\\Primer avance: PostgreSQL y MongoDB}
\author{José, Daniel, Kembly y Fernanda\\
  Ingeniería en Computadores\\Instituto Tecnológico de Costa Rica}
\date{28 de octubre de 2026}

\begin{document}
\maketitle
\begin{abstract}
Este informe documenta el avance de almacenamiento relacional y documental de
SmartCity Analytics. El alcance son las secciones Uso de PostgreSQL y Uso de
MongoDB del enunciado: modelos, 3FN, estructuras, índices, SQL avanzado,
documentos, CRUD y Aggregation Framework, además de veinte consultas.
La generación Madrid v2 acredita 748.889 lecturas por motor y la ejecución
técnica de las consultas. La ejecución PostgreSQL normalizada se contrasta con
la evidencia MongoDB conservada y con los resultados relacionales anteriores.
El modelo relacional separa las dependencias funcionales y reconstruye el
contrato documental mediante vistas, conservando la identidad de las mediciones.
Los datos administrativos sintéticos se distinguen de las observaciones públicas.
El laboratorio académico complementa MG01 y MG07 con lecturas sintéticas por
sensor y tipo de dispositivo, y demuestra recuperación de texto narrativo.
El almacenamiento distribuido MongoDB sigue sin demostración; su requisito
permanece identificado dentro del avance. Hadoop y Spark corresponden a la
continuación del proyecto.
\end{abstract}
\newpage
\tableofcontents
\newpage
\section{Problema y alcance del avance}
Las fuentes urbanas publican datos con diferentes estructuras, frecuencias,
reglas de calidad e identificadores. SmartCity Analytics estudia cómo almacenar
estos datos y consultarlos sin perder su procedencia ni atribuirles información
que no existe en los originales.

El primer avance se concentra en PostgreSQL y MongoDB: modelos, restricciones,
índices, población de datos, diez consultas por motor y evidencia de ejecución.
El plan del equipo sitúa este avance el 28 de octubre y la entrega final el
25 de noviembre de 2026. Hadoop, Spark y la interfaz corresponden al hito final.
La división procede del plan de trabajo; el enunciado general no incluye una
rúbrica independiente del primer avance \cite{plan,enunciado}.

\subsection{Objetivos del avance}
\begin{itemize}
\item Representar catálogos y observaciones con integridad relacional.
\item Representar lecturas y sus metadatos como documentos BSON validados.
\item Demostrar consultas, agregaciones, índices y operaciones CRUD solicitadas.
\item Comprobar replay, conflictos y respaldo/restauración en entornos locales.
\item Documentar qué resultados están acreditados y qué requisitos faltan.
\end{itemize}

\section{Fuentes, preparación y procedencia}
Las fuentes consideradas son calidad del aire, meteorología y tráfico del
Ayuntamiento de Madrid. Sus catálogos oficiales y la investigación del repositorio
documentan licencia CC BY 4.0; deben conservarse atribución, snapshot y cambios
realizados \cite{aire,meteo,trafico,contrato}.

\begin{longtable}{p{0.18\textwidth}p{0.34\textwidth}p{0.38\textwidth}}
\toprule
Fuente & Información & Estado en la generación anual \\
\midrule\endhead
Aire & Observaciones horarias y flags de calidad. &
1.076.813 observaciones aceptadas de 2025. \\
Meteorología & Variables y estaciones meteorológicas. &
Los 69.408 slots auditados de enero permanecen en cuarentena;
el reloj no está acreditado para habilitar su aceptación. \\
Tráfico & Intervalos y variables publicadas por punto de medida. &
677.724 filas seleccionadas de enero permanecen en cuarentena;
la revisión incompatible del contrato sigue pendiente. \\
\bottomrule
\end{longtable}

Solo aire constituye una fuente de mediciones aceptada en la carga anual.
Investigar tres fuentes o disponer de inventarios auxiliares no demuestra por sí
solo el requisito de dos fuentes integradas.

\subsection{Generación local v2: dos fuentes, sin modificar v1}
La implementación local autorizada el 27/10 prepara un paquete nuevo de
\textbf{748.889 lecturas reales de enero de 2025}: 91.959 de aire y 656.930 de
tráfico. Se excluyen 17.774 registros C30 y 3.020 flags vacíos de tráfico, además
de los rechazos de aire. No se suman estos datos a la carga anual v1 ni se
presentan como una ampliación de 1.733.743 registros.

\archivo{sprint1/basico/preparar_v2.py} crea la generación
\archivo{expanded-v2-2025-01} con cuarentena y manifiesto de hashes.
La preparación está terminada; las cargas nativas y las consultas deben
acreditarse con el estado y las salidas de los nuevos informes de ejecución.
Disponer del archivo preparado no demuestra por sí solo una carga en cada motor.

Como complemento separado se preparan 14.388 accidentes de 2026,
deduplicados desde 33.322 filas de personas \cite{accidentes}.
Se eliminan sexo, edad, lesiones, dirección y expediente en claro. Las fechas
se conservan como fechas publicadas sin inventar un reloj UTC; no se unen
estos accidentes de 2026 con las observaciones de enero de 2025.

\subsection{Calidad y conteos}
El snapshot de aire produjo 1.103.232 slots de 2025: 1.076.813 aceptados y
26.419 rechazados. Los rechazos corresponden a 26.167 flags del publicador,
126 horas locales inexistentes y 126 ambiguas sin evidencia para resolverlas.
La ampliación anual contiene enero; no se le suman de nuevo sus 91.959 aceptadas
ni las muestras de prueba \cite{cierre}.

El campo de tráfico \archivo{periodo_integracion} cuenta muestras, no minutos.
La duración documentada es de 15 minutos. Esta corrección está en la propuesta
v2 y no habilita automáticamente la carga de tráfico bajo v1 \cite{v2}.

La versión local v2 aplica esa corrección en una base nueva: duración 15 y
conteo de muestras positivo, incluso cuando el conteo es 30. Solo admite URB/M30
con flag N y tiempo local inequívoco. Los negativos opcionales se marcan como
faltantes; URB no recibe una velocidad inventada y carga es un índice.
Los distritos provienen del inventario disponible, sin vigencia histórica
confirmada para enero de 2025.

\subsection{Identidad y trazabilidad}
Se distinguen estación, canal lógico y sensor físico. Un canal no acredita la
identidad de un dispositivo individual. El contrato anual deriva identificadores
mediante SHA-256 de contenido canónico y conserva fuente, versión del archivo,
referencia de fila/columna, calidad, tiempo original y transformación.

\archivo{source_version} identifica el snapshot y debe coincidir con
\archivo{raw_reference.sha256}. Las horas de aire se interpretan en
\archivo{Europe/Madrid}; las ambigüedades no se resuelven inventando un fold.
Los conflictos de contenido no deben sobrescribir observaciones aceptadas.

\subsection{Datos administrativos sintéticos}
La evidencia documental de la respuesta docente del 10/10 permite atributos
administrativos aleatorios. El paquete anual registra una organización,
5.133 usuarios y 5.133 sensores sintéticos deterministas, con semilla
\archivo{20261010}, versión \archivo{admin-demo-v2} y procedencia explícita.
Los correos usan \archivo{example.invalid}. Esto no acredita inventario físico
ni autoriza mediciones inventadas \cite{permiso,contrato}.

Los sensores demo tienen identificadores propios, pero no se enlazan a canales
o lecturas observadas: la referencia \archivo{sensor_id} en los canales y en
la proyección documental de las lecturas reales permanece nula.
La autorización se conoce por el registro
aportado al repositorio, no por una captura original adjunta a este informe.

\section{Modelo relacional y documental}
\subsection{Estructura relacional normalizada}
El esquema de entrega se define en
\archivo{sprint1/postgres/schema_v2_3nf.sql} \cite{pg3nf}. Conserva las doce entidades del
núcleo anterior y añade dos relaciones auxiliares: tipos de dispositivo y
asignaciones sensor--canal. El modelo definitivo contiene catorce tablas.
El esquema de práctica \archivo{sprint1}, que usa \archivo{readings} y
\archivo{storage_id}, es independiente y no representa este modelo.

\begin{longtable}{>{\raggedright\arraybackslash}p{0.25\textwidth}>{\raggedright\arraybackslash}p{0.26\textwidth}>{\raggedright\arraybackslash}p{0.40\textwidth}}
\toprule Tabla & Claves candidatas & Función \\
\midrule\endhead
sources & source\_id & Dataset y publicador. \\
source\_versions & source\_version & Snapshot y fuente. \\
zones & zone\_id & Zona administrativa y distrito padre opcional. \\
locations & location\_id & Nombre del sitio, coordenadas y zona opcional. \\
stations & station\_id; fuente/código & Sitio lógico del inventario. \\
sensor\_categories & category\_id & Clasificación administrativa de dispositivos. \\
measurement\_types & measurement\_type; fuente/código & Magnitud medida y unidad; independiente del tipo físico. \\
organizations & organization\_id & Organismo y responsable. \\
users & user\_id; email & Usuario de aplicación. \\
sensor\_types & sensor\_type & Tipo de dispositivo y categoría. \\
sensors & sensor\_id & Dispositivo administrativo, instalación y estado. \\
channels & channel\_id; estación/clave & Serie lógica de una magnitud. \\
channel\_sensors & channel\_id & Asignación conocida de canal a dispositivo. \\
observations & measurement\_id; canal/UTC & Hecho observado y procedencia variable. \\
\bottomrule
\end{longtable}

\clearpage
\begin{figure}[p]
\centering
\begin{tikzpicture}[x=3.65cm,y=2.1cm,
 entidad/.style={draw,fill=white,rounded corners,align=center,font=\scriptsize,minimum width=2.9cm,minimum height=0.75cm},
 relacion/.style={-{Latex},thin}]
\node[entidad] (src) at (0,0) {sources};
\node[entidad] (ver) at (1,0) {source\_versions};
\node[entidad] (org) at (2,0) {organizations};
\node[entidad] (usr) at (3,0) {users};
\node[entidad] (sta) at (0,-1) {stations};
\node[entidad] (mag) at (1,-1) {measurement\_types};
\node[entidad] (cat) at (2,-1) {sensor\_categories};
\node[entidad] (zon) at (3,-1) {zones};
\node[entidad] (cha) at (0,-2) {channels};
\node[entidad] (sen) at (1,-2) {sensors};
\node[entidad] (typ) at (2,-2) {sensor\_types};
\node[entidad] (loc) at (3,-2) {locations};
\node[entidad] (obs) at (0,-3) {observations};
\node[entidad] (lnk) at (1,-3) {channel\_sensors};
\draw[relacion] (ver) -- (src);
\draw[relacion] (sta) -- (src);
\draw[relacion] (sta.south) -- ++(0,-0.22) -| ([xshift=-3mm]loc.west) -- (loc.west);
\draw[relacion] (mag) -- (src);
\draw[relacion] (typ) -- (cat);
\draw[relacion] (sen) -- (typ);
\draw[relacion] (sen) -- (sta);
\draw[relacion] (sen.north) -- (1,-1.5) -- (2.5,-1.5)
 -- (2.5,-0.5) -- (2,-0.5) -- (org.south);
\draw[relacion] (cha) -- (sta);
\draw[relacion] (cha) -- (mag);
\draw[relacion] (lnk) -- (cha);
\draw[relacion] (lnk) -- (sen);
\draw[relacion] (obs) -- (cha);
\draw[relacion] (obs.east) -- ++(0.10,0) |- ([yshift=-4mm]ver.south) -- (ver.south);
\draw[relacion] (loc) -- (zon);
\draw[relacion] (loc.east) -- ++(0.15,0) -- ++(0,2.5) -| (ver.north);
\draw[relacion] (zon) edge[loop right] (zon);
\end{tikzpicture}
\caption{Modelo relacional completo de catorce tablas. Cada flecha parte de
la relación hija y apunta a la relación referenciada. Usuarios es independiente.}
\end{figure}
\clearpage

El E--R con atributos y cardinalidades se conserva en
\archivo{docs/arquitectura/diagrama_3nf.md}.

Las relaciones admiten cero o muchos hijos, salvo canal--asignación, limitada
a una fila por canal. Zona, distrito padre, versión de inventario, ubicación de
estación, tipo de dispositivo y organismo son referencias opcionales; las demás
FK exigen un padre por hijo. Una zona local puede pertenecer a un distrito.
\archivo{location_geography} reconstruye ambos nombres, sin repetirlos en cada
ubicación. Los 21 nombres de distrito publicados no acreditan vigencia histórica.

\subsection{Lectura del modelo y cardinalidades}
Una estación identifica un punto lógico de la fuente de datos. Un canal identifica
la serie de una magnitud en esa estación, y una observación registra su valor en
un instante. No se equipara un canal con un dispositivo físico: esa asociación
se almacena por separado en \archivo{channel_sensors}, únicamente cuando se conoce.
En la generación observada actual no hay asignaciones físicas acreditadas.

La tabla siguiente complementa las flechas del diagrama. Cada fila describe la
referencia desde la tabla hija hacia su padre. Una referencia obligatoria exige
exactamente un padre; una opcional permite cero o uno. Cada padre admite cero o
muchas filas hijas, salvo la asignación de canal indicada al pie.

\begin{longtable}{>{\raggedright\arraybackslash}p{0.25\textwidth}>{\raggedright\arraybackslash}p{0.25\textwidth}>{\raggedright\arraybackslash}p{0.25\textwidth}>{\raggedright\arraybackslash}p{0.14\textwidth}}
\toprule Tabla hija & Clave foránea & Tabla padre & Referencia \\
\midrule\endhead
\archivo{source_versions} & \archivo{source_id} & \archivo{sources} & Obligatoria \\
\archivo{zones} & \archivo{parent_zone_id} & \archivo{zones} & Opcional \\
\archivo{locations} & \archivo{zone_id} & \archivo{zones} & Opcional \\
\archivo{locations} & \archivo{inventory_version} & \archivo{source_versions} & Opcional \\
\archivo{stations} & \archivo{source_id} & \archivo{sources} & Obligatoria \\
\archivo{stations} & \archivo{location_id} & \archivo{locations} & Opcional \\
\archivo{measurement_types} & \archivo{source_id} & \archivo{sources} & Obligatoria \\
\archivo{sensor_types} & \archivo{category_id} & \archivo{sensor_categories} & Obligatoria \\
\archivo{sensors} & \archivo{sensor_type} & \archivo{sensor_types} & Opcional \\
\archivo{sensors} & \archivo{station_id} & \archivo{stations} & Obligatoria \\
\archivo{sensors} & \archivo{organization_id} & \archivo{organizations} & Opcional \\
\archivo{channels} & \archivo{station_id} & \archivo{stations} & Obligatoria \\
\archivo{channels} & \archivo{measurement_type} & \archivo{measurement_types} & Obligatoria \\
\archivo{channel_sensors} & \archivo{channel_id} & \archivo{channels} & Obligatoria \\
\archivo{channel_sensors} & \archivo{sensor_id} & \archivo{sensors} & Obligatoria \\
\archivo{observations} & \archivo{channel_id} & \archivo{channels} & Obligatoria \\
\archivo{observations} & \archivo{source_version} & \archivo{source_versions} & Obligatoria \\
\bottomrule
\end{longtable}

\archivo{channel_sensors.channel_id} es también clave primaria: un canal admite
cero o una asignación, mientras que un sensor puede estar asociado a varios
canales. \archivo{users} no tiene claves foráneas en este modelo; no se dibujan
relaciones con usuarios que el esquema no implementa. La autorreferencia de
\archivo{zones} permite un distrito padre opcional para una zona local; el trigger
comprueba el nivel del padre.

\subsection{Dependencias funcionales y tercera forma normal}
La condición de 3FN exige que, en cada DF no trivial \(X\to A\), \(X\) sea
superclave o \(A\) atributo primo. La cobertura de DF y las claves candidatas se
registran en \archivo{sprint1/postgres/dependencias_3nf.json}. Se calcularon los
cierres de todos los subconjuntos de atributos y se comparó el inventario con
las columnas reales de PostgreSQL. Con esa cobertura, las catorce relaciones
satisfacen 3FN y BCNF.

Estación--fuente y magnitud--fuente permanecen en sus catálogos; canales y
observaciones no repiten fuente. La asignación sensor--canal evita duplicar la
estación del dispositivo. Tipo--categoría se conserva en \archivo{sensor_types}.

La vista contractual calcula la hora local desde UTC, proyecta constantes y
obtiene calidad, estado temporal, rol e intervalo desde la política por fuente.
También repone el SHA del snapshot en la referencia original. La reconstrucción
conserva el DTO v2, sus identificadores y valores.

La reunión por claves de catálogo conserva la información; las asignaciones
opcionales usan LEFT JOIN. Los triggers comprueban fuentes, estación del sensor,
inmutabilidad y aceptación v2. Se acreditan ocho rechazos de operaciones, sin
aislar cada CHECK o FK; las pruebas positivas y PG07 utilizan rollback.

JSONB es un dominio SQL tipado indivisible; no se afirma normalización escalar
recursiva. Los NULL opcionales denotan ausencia y no se usan como claves. No se
infieren DF de coincidencias del lote o de condiciones solo para valores no nulos.

\subsection{Dispositivos, magnitudes y procedencia}
\archivo{sensor_types} representa tipos de dispositivos administrativos;
\archivo{measurement_types} representa variables y unidades medidas. No hay FK
entre magnitudes y categorías físicas. Los doce tipos demostrativos y sus
categorías son explícitamente sintéticos. Se conservan 5.133 sensores y usuarios
demo y un organismo demo, sin enlaces inventados con canales observados.

\subsection{Modelo documental y compatibilidad}
MongoDB conserva \archivo{readings}, \archivo{incidents}, \archivo{derived_events},
\archivo{alerts}, \archivo{organizations}, \archivo{users} y \archivo{sensors}.
Las lecturas reúnen valor, calidad, reloj y procedencia, con BSON Date e
\archivo{_id=measurement_id}. Los objetos \archivo{measurements} y
\archivo{provenance} representan información semiestructurada. Esa flexibilidad
y los patrones de acceso justifican el modelo documental. La recuperación de
información narrativa se demuestra por separado con documentos sintéticos, sin
atribuir contenido no estructurado a toda lectura por el solo hecho de usar BSON.

\archivo{observation_contract} conserva la interfaz compartida, identidad y
payload, sin cambiar MongoDB ni establecer FK o transacciones distribuidas entre
motores. Persisten los límites de sensores físicos, geografía y reglas demo.

\clearpage
\subsection{Diagrama documental de la generación Madrid v2}
\begin{figure}[ht]
\centering
\begin{tikzpicture}[x=1cm,y=1cm,
  doc/.style={draw,rounded corners,align=left,text width=4.55cm,
  inner sep=6pt,font=\footnotesize,anchor=north},
  ref/.style={-{Latex[length=2mm]},thick},
  lab/.style={font=\scriptsize,fill=white,inner sep=2pt}]
\node[doc] (r) at (0,0) {\textbf{readings} [V]\\
  \texttt{\_id = measurement\_id}: string\\
  \texttt{channel\_id, station\_id}: string\\
  \texttt{source\_id, source\_version}: string\\
  \texttt{observed\_at}: BSON Date\\
  \texttt{value}: número; \texttt{unit}: string/null\\
  \texttt{sensor\_id}: null; origen observado\\
  \texttt{raw\_reference, provenance}: objetos\\
  \texttt{measurements}: objeto opcional\\
  Calidad, zona y versión contractual.};
\node[doc] (e) at (5.4,0) {\textbf{derived\_events} [V]\\
  \texttt{\_id, measurement\_id}: string\\
  \texttt{rule\_version}: string\\
  \texttt{threshold, value}: números\\
  \texttt{unit}: string; \texttt{zone\_id}: string/null\\
  \texttt{observed\_at, generated\_at}: Date\\
  \texttt{data\_origin}: derived\\
  \texttt{official\_event}: false\\
  Derivación por regla de demostración.};
\node[doc] (a) at (10.8,0) {\textbf{alerts} [V]\\
  \texttt{\_id, event\_id}: string\\
  \texttt{zone\_id, rule\_version}: string\\
  \texttt{observed\_at, emitted\_at}: Date\\
  \texttt{data\_origin}: derived\\
  \texttt{replay}: true\\
  \texttt{delivery}: local\_collection\_only\\
  \texttt{official\_alert}: false\\
  No constituye notificación externa.};
\draw[ref] (e.west) -- (r.east);
\draw[ref] (a.west) -- (e.east);
\node[doc] (o) at (0,-7) {\textbf{organizations} [sin V]\\
  \texttt{\_id = organization\_id}: string\\
  \texttt{name}: string\\
  \texttt{responsible\_person\_or\_role}: string\\
  \texttt{evidence\_reference}: string\\
  \texttt{provenance}: objeto\\
  \texttt{data\_origin}: synthetic.};
\node[doc] (s) at (5.4,-7) {\textbf{sensors} [sin V]\\
  \texttt{\_id = sensor\_id}: string\\
  Nombre, categoría y estación: string\\
  \texttt{organization\_id}: string\\
  Instalación y fecha de estado: string\\
  \texttt{provenance}: objeto\\
  Origen sintético; sin \texttt{sensor\_type}.};
\node[doc] (u) at (10.8,-7) {\textbf{users} [sin V]\\
  \texttt{\_id = user\_id}: string\\
  \texttt{name, email}: string\\
  \texttt{consent\_reference}: string\\
  \texttt{demo\_only}: boolean\\
  \texttt{provenance}: objeto\\
  Sin vínculo usuario--sensor definido.};
\draw[ref] (s.west) -- (o.east);
\node[doc] (i) at (0,-12) {\textbf{incidents} [V]\\
  \texttt{\_id, category}: string\\
  \texttt{reported\_date}: string (2026)\\
  \texttt{zone\_id}: string/null\\
  Fuente y versión: string\\
  \texttt{raw\_reference}: objeto\\
  \texttt{person\_fields\_removed}: true.};
\node[align=left,text width=9.8cm,anchor=north west,font=\footnotesize] at (2.7,-12) {
  \textbf{Leyenda.} Las flechas representan referencias lógicas del contenido,
  no claves foráneas MongoDB. [V] indica validador existente; los tipos mostrados
  describen los documentos cargados, no una restricción de tipo sobre cada campo.
  Las restricciones efectivas se detallan en el cargador y el catálogo nativo.\\[4pt]
  Estación, canal, zona y categoría remiten a catálogos externos a estas siete
  colecciones. No hay vínculo entre sensores administrativos y lecturas públicas:
  \texttt{sensor\_id} es nulo. Los incidentes no se unen temporalmente a las lecturas.};
\end{tikzpicture}
\caption{Colecciones oficiales y referencias lógicas de Madrid v2.}
\end{figure}

\begin{longtable}{>{\raggedright\arraybackslash}p{0.23\textwidth}>{\raggedright\arraybackslash}p{0.67\textwidth}}
\toprule Colección & Validación e índices, además de \texttt{\_id} \\
\midrule\endhead
readings & strict/error; UNIQUE (channel\_id, observed\_at, timestamp\_role);
(source\_id, measurement\_type, observed\_at); (zone\_id, measurement\_type). \\
derived\_events & error; (rule\_version, observed\_at). \\
alerts & error; (zone\_id, emitted\_at). \\
incidents & error; (category). \\
Administrativas & Solo el índice \texttt{\_id}; sin validadores adicionales. \\
\bottomrule
\end{longtable}
\clearpage


\clearpage
\section{Correspondencia con las entidades exigidas}
Los nombres de columna se expresan en inglés; su correspondencia con el
vocabulario del enunciado se establece a continuación. La normalización permite
obtener atributos mediante relaciones, sin duplicarlos en una tabla.

\begingroup
\small
\setlength{\tabcolsep}{4pt}
\begin{longtable}{>{\raggedright\arraybackslash}p{0.24\textwidth}>{\raggedright\arraybackslash}p{0.44\textwidth}>{\raggedright\arraybackslash}p{0.23\textwidth}}
\toprule Atributo exigido & Implementación & Nulabilidad y origen \\
\midrule\endhead
Sensor: id\_sensor & \archivo{sensors.sensor_id}, PK & No; demo sintético. \\
Sensor: nombre & \archivo{sensors.name} & No; sintético. \\
Sensor: tipo\_sensor & \archivo{sensors.sensor_type} $\to$ \archivo{sensor_types.name} & FK opcional; los tipos cargados son sintéticos. \\
Sensor: fecha\_instalacion & \archivo{sensors.installation_date} & Sí; fechas demo. \\
Sensor: estado & \archivo{sensors.administrative_status} & Sí; active/inactive demo. \\
Ubicación: id\_ubicacion & \archivo{locations.location_id}, PK & No; identidad derivada. \\
Ubicación: nombre\_zona & \archivo{locations.zone_id} $\to$ \archivo{zones.name} & Relación opcional; nombre administrativo publicado. \\
Ubicación: distrito & \archivo{location_geography.district_id}/district\_name, derivados de zones y su padre & Opcional; 55 ubicaciones sin asignación. \\
Ubicación: coordenadas & \archivo{locations.latitude}/longitude y coordinate\_crs & Coordenadas opcionales por pares; presentes en 5.133 filas. CRS sin acreditación definitiva. \\
Categoría: id\_categoria & \archivo{sensor_categories.category_id}, PK & No; clasificación demo. \\
Categoría: nombre\_categoria & \archivo{sensor_categories.name} & No; etiqueta sintética explícita. \\
Usuario: id\_usuario & \archivo{users.user_id}, PK & No; sintético. \\
Usuario: nombre & \archivo{users.name} & No; sintético. \\
Usuario: correo & \archivo{users.email}, UNIQUE & No; dominio \archivo{example.invalid}. \\
Organismo: id\_organismo & \archivo{organizations.organization_id}, PK & No; sintético. \\
Organismo: nombre & \archivo{organizations.name} & No; sintético. \\
Organismo: responsable & \archivo{organizations.responsible_person_or_role} & Sí; rol ficticio de demostración. \\
\bottomrule
\end{longtable}
\endgroup

Los 17 atributos tienen representación directa o relacional. El enunciado no
establece NOT NULL para todos ellos; las ausencias se conservan sin sustituirlas
por datos reales inventados. El nombre propio del sitio es
\archivo{locations.name}, distinto del nombre de zona; existen 22 nombres de sitio
ausentes. Las coordenadas del inventario y los nombres administrativos no prueban
la ubicación de cada dispositivo en enero de 2025.

El catálogo mantiene tres niveles territoriales: ciudad, distrito y zona local.
La FK de una zona local solo admite un padre de nivel distrito. No se crean
subzonas observadas ni se asignan distritos por proximidad. Las pruebas de esa
relación utilizan exclusivamente fixtures identificadas y revertidas.

\section{Consultas obligatorias y cobertura}
Las siguientes tablas resumen la evidencia anual, las pruebas controladas y la
ruta local v2, diferenciadas por fila. Una implementación nueva se acredita con
su informe de ejecución, no solo con que el archivo exista.
``Demo'' significa entidades
administrativas sintéticas autorizadas; no inventario físico observado.
``Registrada'' no equivale a aceptación académica ni reproducibilidad actual.

\subsection{PostgreSQL}
\begin{longtable}{p{0.08\textwidth}p{0.39\textwidth}p{0.43\textwidth}}
\toprule ID & Consulta exigida & Evidencia y límite \\
\midrule\endhead
PG01 & Sensores por tipo. & Registrada sobre sensores demo. \\
PG02 & Sensores por ubicación. & Registrada sobre sensores demo y catálogo. \\
PG03 & Cantidad de activos e inactivos. & Registrada con estados administrativos demo. \\
PG04 & Sensores por organismo responsable. & Registrada con organismos asignados a sensores demo. \\
PG05 & Ubicaciones con más sensores. & Ranking registrado sobre sensores demo. \\
PG06 & Sensores instalados por año. & Registrada con fechas sintéticas. \\
PG07 & Sensores sin actividad registrada. & Funcional en prueba controlada; bloqueada en el lote observado por falta de relación sensor--canal. \\
PG08 & Ranking de zonas monitoreadas. & Registrada sobre sensores demo; revisar alcance geográfico. \\
PG09 & Consulta utilizando vistas. & Vista anual de actividad por canal. \\
PG10 & Procedimiento o función. & Función anual que cuenta canales de una estación. \\
\bottomrule
\end{longtable}

La regresión de integración fue corregida: el archivo
\archivo{sprint1/postgres/queries.sql} vuelve a contener PG01--PG10 para
\archivo{smartcity}; el núcleo local usa \archivo{staging_queries.sql}.
El parser anual reconoce las diez consultas y tiene pruebas de regresión.
Una nueva prueba PostgreSQL ejecuta el archivo canónico y contrasta las filas
de cada consulta con resultados esperados sobre fixtures revertidas al finalizar.

PG07 requiere una relación conocida sensor--canal antes de evaluar ausencia
de lecturas en enero. Excluye relaciones desconocidas y usa una ventana
semiabierta en hora de Madrid. La prueba incluye ambos límites, un sensor con
canal sin lecturas y un sensor sin canal; no acredita identidad física observada.

El núcleo demuestra SELECT, filtros, ordenamiento, DISTINCT, JOIN, GROUP BY,
HAVING, agregaciones, CTE, ventanas, subconsultas y fechas en distintos archivos.
Esa variedad técnica no reemplaza las preguntas obligatorias sobre sensores.

\subsection{MongoDB}
\begin{longtable}{p{0.08\textwidth}p{0.39\textwidth}p{0.43\textwidth}}
\toprule ID & Consulta exigida & Evidencia y límite \\
\midrule\endhead
MG01 & Lecturas de un sensor específico. & Búsqueda registrada por canal lógico, no dispositivo físico. \\
MG02 & Lecturas en un rango de fechas. & V2: rango BSON semiabierto, contrastado con conteo independiente del archivo. \\
MG03 & Máximos y mínimos. & V2: extremos de aire por variable/unidad, comparados con Python. \\
MG04 & Eventos críticos detectados. & V2: regla demo de intensidad superior a 2.000 veh/h; eventos derivados, no oficiales. \\
MG05 & Alertas emitidas por zona. & V2: emisión actual de replay en colección local; zonas del inventario disponible. \\
MG06 & Promedio de mediciones por zona. & V2: solo zonas conocidas; vigencia histórica no confirmada. \\
MG07 & Promedio por tipo de sensor. & V2: por fuente/variable/unidad; alternativa por tipo lógico, no físico. \\
MG08 & Incidentes por categoría. & V2: accidentes reales de 2026 deduplicados y anonimizados. \\
MG09 & Consulta con agregaciones. & V2: conteos por fuente y flag de calidad. \\
MG10 & Consulta con filtros múltiples. & V2: fecha, fuente, flag, valor y zona; resultado comparado con Python. \\
\bottomrule
\end{longtable}

El harness anual contiene algunas afirmaciones constantes en
\archivo{query_evidence}; su verificador controla estructura y procedencia,
no ejecuta cada pipeline. Por ello, no se declaran siete consultas anuales
completamente acreditadas solo por esas banderas. Agrupar una zona nula tampoco
demuestra un análisis por zona geográfica.

La ruta nueva no reutiliza esas banderas como prueba. El archivo
\archivo{sprint1/basico/consultas_mongodb_v2.js} contiene los diez pipelines y
filtros. \archivo{scripts/validar_v2.py} compara cantidades, extremos y promedios
con cálculos independientes sobre el JSONL verificado. Los informes nuevos
conservan parámetros, muestras y resultados agregados de cada consulta.

MG04 usa una regla académica de demostración, no un límite legal. MG05 guarda
\archivo{emitted_at} actual, \archivo{replay=true} y
\archivo{delivery=local_collection_only}; no fabrica alertas oficiales pasadas
ni envía notificaciones externas. El uso de distritos del inventario se declara
como \archivo{current_inventory_not_historically_verified}.

Las operaciones insertOne, insertMany, find, findOne, updateOne, updateMany y
deleteOne se prueban en la integración con colecciones de demostración aisladas.
Las actualizaciones y borrados no deben aplicarse a las observaciones originales.

\subsection{Ficha mínima para completar cada consulta}
Daniel y Kembly deben registrar, para cada ID:
\begin{enumerate}
\item Pregunta, contrato, esquema/colección y generación utilizada.
\item Archivo y SQL o pipeline exacto, con sus parámetros.
\item Resultado obtenido y comprobación independiente del valor esperado.
\item Índice/plan cuando corresponda, fecha, entorno y responsable de revisión.
\item Origen observado, derivado o sintético y limitaciones de interpretación.
\end{enumerate}
La matriz editable del repositorio está en
\archivo{docs/equipo/MATRIZ_CONSULTAS_SPRINT1.md}.

\section{Validación y resultados disponibles}
\subsection{Demostración local del 27 de octubre}
Se ejecutaron 32 pruebas Python y 8 Node, todas aprobadas. La integración
PostgreSQL 16.6/MongoDB 7.0.16 terminó en PASS con tres ejemplos publicados,
uno de aire, uno meteorológico y uno de tráfico, adaptados en copias de prueba.
Los ejemplos de meteo/tráfico en staging no son registros aceptados de la carga anual.

La integración comprobó restricciones SQL/BSON, índices, consultas, replay,
conflictos en cuarentena, CRUD aislado y respaldo/restauración lógico de ambos
motores. También probó ocho conexiones PostgreSQL sobre una clave y dos payloads:
una aceptación, tres duplicados y cuatro cuarentenas. No es estrés sostenido ni
prueba de concurrencia MongoDB entre procesos \cite{integracion,respaldos}.

Los contenedores fueron locales, sin puertos publicados y con red deshabilitada.
Se eliminaron al finalizar. El informe generado fue
\archivo{/tmp/opencode/smartcity-s1-integration-t7d2ibo2.json}; es un artefacto
local, no un archivo disponible en cualquier clon.

\subsection{Carga anual: evidencia previa, no reejecución de hoy}
El mínimo es de 100.000 registros PostgreSQL y 500.000 documentos MongoDB
según el enunciado y el plan del equipo \cite{enunciado,plan}.
\begin{longtable}{p{0.48\textwidth}p{0.2\textwidth}p{0.2\textwidth}}
\toprule Control registrado & PostgreSQL & MongoDB \\
\midrule\endhead
Lecturas/documentos aceptados & 1.076.813 & 1.076.813 \\
Metadatos y valores comparados & 1.076.813 & 1.076.813 \\
Replay del lote completo & Sin aumento & Sin aumento \\
Mediciones sintéticas & 0 & 0 \\
Mínimo de volumen del motor & Superado localmente & Superado localmente \\
Segunda fuente de mediciones aceptada & No & No \\
\bottomrule
\end{longtable}

Estos resultados proceden de los manifiestos
\archivo{research/datasets/results/expanded-2025/postgresql_full.json} y
\archivo{mongodb_full.json}, y del cierre documentado \cite{cierre}.
La revisión actual contrastó los hashes de sus cuatro scripts registrados con
el código disponible; no repitió la carga masiva ni verificó todos los insumos privados.

Los mínimos locales no prueban que el servidor compartido tenga esos datos.
Tampoco constituyen una comparación de rendimiento entre motores: almacenamiento,
configuración y operaciones deben ser equivalentes antes de comparar tiempos.

\subsection{Reproducción del núcleo}
Desde la raíz del repositorio, con Python 3.10+ y Node.js 18+:
\begin{lstlisting}
bash scripts/validate_sprint1.sh
\end{lstlisting}
Para integrar motores se requiere Docker local y las imágenes ya descargadas
\archivo{postgres:16.6-alpine} y \archivo{mongo:7.0.16}:
\begin{lstlisting}
PYTHONDONTWRITEBYTECODE=1 python3 scripts/test_sprint1_engines.py
\end{lstlisting}
El harness no descarga imágenes. Las instrucciones completas y su alcance están
en \archivo{ENTREGA_SPRINT1.md}. Los clientes básicos psycopg/PyMongo aún
necesitan prueba de extremo a extremo con autenticación real.

\subsection{Respaldo y restauración}
La prueba usa pg\_dump/pg\_restore y mongodump/mongorestore, restaura en otra
base desechable y compara contenido y estructura. Esto acredita el caso
controlado, no una política de respaldo productiva, recuperación tras caída,
retención externa o restauración del servidor compartido.

\subsection{Qué evidencia falta adjuntar}
\begin{itemize}
\item Salidas verificables de cada una de las veinte consultas en la ruta elegida.
\item Diagrama E--R completo y dependencias funcionales revisados por Daniel.
\item Modelo documental final, ejemplos BSON y revisión de Kembly.
\item Capturas legibles y registro de revisión cruzada, sin credenciales ni hosts privados.
\item Comprobación de acceso, conteos y respaldo compartidos si se usan en la entrega.
\end{itemize}

\section{Dictamen y limitaciones}
El modelo PostgreSQL representa las cinco entidades mínimas y sus 17 atributos,
mediante columnas directas o relaciones normalizadas. Las dependencias de
fuente, reloj, contrato, tipo y categoría se separan del hecho observado.
Las catorce relaciones satisfacen 3FN respecto de la cobertura funcional
documentada y los dominios SQL tipados. Las dos tablas auxiliares son parte del
modelo definitivo y se incluyen en el E--R completo.

Los tipos administrativos y categorías demo no representan aparatos acreditados
del inventario municipal. Las fechas, estados, responsables y usuarios
sintéticos permanecen identificados. Los vínculos físicos desconocidos no se
completan por inferencia; PG07 conserva esa distinción. Las coordenadas y
asignaciones territoriales no acreditan vigencia histórica individual.

MongoDB conserva su contrato y sus evidencias nativas: lecturas reales, eventos
por regla académica, alertas de replay local e incidentes anonimizados de 2026.
Las consultas por canal o variable no se presentan como identificación de un
dispositivo físico. No se atribuyen alertas oficiales ni causalidad entre períodos.
El laboratorio separado demuestra MG01 por identidad de sensor ficticio y MG07
mediante tipos administrativos, sin equiparar canal y dispositivo. Tres documentos
con texto libre, atributos heterogéneos y revisiones demuestran almacenamiento y
recuperación de información narrativa sintética. La procedencia limita la
interpretación: no acredita dispositivos físicos ni reportes ciudadanos reales.
El almacenamiento distribuido MongoDB permanece sin demostración. Por ese
requisito abierto no se declara el Sprint 1 académicamente completo.

\section{Conclusiones}
La separación entre almacenamiento relacional y vista contractual permite
normalizar PostgreSQL y conservar la identidad compartida con MongoDB. Las
restricciones interrelacionales se mantienen mediante FK y triggers sin repetir
atributos determinados por catálogos. La verificación combina prueba formal de
DF, ejecución nativa, comparación independiente y reconciliación de identidades.

\begin{samepage}
Los mínimos de volumen se evalúan por motor y por generación, sin sumar replay
ni datos administrativos. La ejecución local no acredita el estado de servidores
compartidos. Hadoop/HDFS, Spark, dashboard completo, análisis distribuidos y la
comparación final entre cuatro motores corresponden a la continuación de
noviembre. Esta exclusión no se extiende al almacenamiento distribuido mencionado
en MongoDB: su demostración y modalidad suficiente permanecen abiertas. Las
vistas materializadas y el particionamiento requieren aclaración de alcance.

\end{samepage}

\begin{thebibliography}{99}
\bibitem{enunciado} Instituto Tecnológico de Costa Rica.
\emph{Enunciado Proyecto SmartCity Analytics, II Semestre 2026}.
PDF facilitado al equipo; entrega general del 25/11/2026.
\bibitem{aire} Ayuntamiento de Madrid.
\emph{Calidad del aire: datos horarios}.
\url{https://datos.madrid.es/dataset/201200-0-calidad-aire-horario}.
\bibitem{meteo} Ayuntamiento de Madrid.
\emph{Datos meteorológicos horarios}.
\url{https://datos.madrid.es/dataset/300352-0-meteorologicos-horarios}.
\bibitem{trafico} Ayuntamiento de Madrid.
\emph{Histórico de datos del tráfico}.
\url{https://datos.madrid.es/dataset/208627-0-transporte-ptomedida-historico}.
\bibitem{accidentes} Ayuntamiento de Madrid, Policía Municipal.
\emph{Accidentes de tráfico con detalle}.
\url{https://datos.madrid.es/dataset/300228-0-accidentes-trafico-detalle}.
Snapshot verificado de 2026; datos personales excluidos de la exportación.
\bibitem{contrato} Repositorio SmartCity.
\emph{Contrato implementado de la generación histórica madrid-draft-v1}.
\archivo{research/datasets/scripts/phase5_postgresql.sql};
\archivo{research/datasets/scripts/unified_synthetic.py};
\archivo{docs/equipo/CONTRATO_UNIFICADO.md}, antecedente de esa generación.
\bibitem{permiso} Repositorio SmartCity.
\emph{Registro de comunicación docente del 10/10/2026:
alcance de los datos administrativos sintéticos}.
\archivo{docs/equipo/EVIDENCIA_AUTORIZACION_DATOS_SINTETICOS.md}.
\bibitem{v2} Repositorio SmartCity.
\emph{Implementación y evidencia de la interfaz Madrid v2}.
\archivo{sprint1/postgres/schema_v2_3nf.sql}, vista
\archivo{observation_contract}; adaptación histórica
\archivo{sprint1/postgres/contract_v2.sql};
\archivo{sprint1/basico/preparar_v2.py};
\archivo{research/datasets/results/expanded-v2-2025-01/execution-d26e40706dcd.json}.
\bibitem{cierre} Repositorio SmartCity.
\emph{Validación histórica de la generación anual madrid-draft-v1}.
\archivo{docs/equipo/CIERRE_VALIDACION_SPRINT1.md}.
\bibitem{alcance} Equipo SmartCity.
\emph{Alcance del Sprint 1: registro de la aclaración docente
del 28/10/2026 comunicada al equipo}.
\archivo{docs/equipo/ALCANCE_CONFIRMADO_SPRINT1.md}.
\bibitem{mongocierre} Equipo SmartCity.
\emph{Cierre técnico nativo MongoDB Madrid v2 y resultados verificables}.
\archivo{sprint1/mongodb/evidencia_final_mongodb.md};
\archivo{sprint1/mongodb/evidencias/20261028T055755-3fea400d/comprobacion_cierre.json};
código de la ejecución: \archivo{804d6a551dc84e22d901b64554c30a712b951f77}.
\bibitem{pg3nf} Equipo SmartCity.
\emph{Modelo PostgreSQL normalizado: dependencias funcionales y validación nativa}.
\archivo{sprint1/postgres/impacto_3nf.md};
\archivo{sprint1/postgres/evidencias/3nf/20261028T071812-55791345.json}.
\bibitem{laboratorio} Equipo SmartCity.
\emph{Demostración académica aislada MongoDB: consultas por sensor, tipos y textos narrativos}.
\archivo{sprint1/mongodb/academico/README.md}; ejecución
\archivo{sprint1/mongodb/evidencias/academico/20261028T161057-eff98453/execution.json}.
\end{thebibliography}

\end{document}
